\subsubsection{Итерация [[K]]}

\begin{enumerate}
  \item Найдём $\alpha^{([[K]])}_1$ для направления $e_1=(1,0)^T$ как вершину параболы:
  \[
  \begin{gathered}
    \alpha^{([[K]])}_1=\textrm{arg}\,\min_{\alpha}f(x^{([[K_PREV]])}+\alpha e_1)=-\dfrac{2ax^{([[K_PREV]])}_1+cx^{([[K_PREV]])}_2+d}{2a}=\\
    =-\dfrac{2\cdot [[A]]\cdot [[X0]]+[[C]]\cdot [[Y0]]+[[D]]}{2\cdot [[A]]}=[[ALPHA_1]]
  \end{gathered}
  \]

  Имеем:
  \[
    \tilde{x}_1^{(1)}=x^{(0)}+\alpha_1^{(1)}e_1=([[X1]],[[Y1]])^T
  \]

  \item Аналогично найдём $\alpha^{([[K]])}_2$ для направления $e_2=(0,1)^T$:
  \[
  \begin{gathered}
    \alpha^{([[K]])}_2=\textrm{arg}\,\min_{\alpha}f(\tilde{x}^{(1)}_1+\alpha e_2)=-\dfrac{2b\tilde{x}^{([[K]])}_2+c\tilde{x}^{([[K]])}_1+e}{2b}=\\
    =-\dfrac{2\cdot [[B]]\cdot [[Y1]]+[[C]]\cdot [[X1]]+[[E]]}{2\cdot [[B]]}=[[ALPHA_2]]
  \end{gathered}
  \]

  Итого имеем:
  \[
    x^{(1)}=\tilde{x}_2^{(1)}=\tilde{x}^{(1)}_1+\alpha_2^{(1)}e_2=([[X2]],[[Y2]])^T
  \]

  \item Графическая интерпретация:
  \begin{figure}[H]
  \centering
  \begin{tikzpicture}
    \begin{axis}[
      width=10cm, height=10cm,
      axis lines=middle,
      xlabel={$x$},
      ylabel={$y$},
      xlabel style={right},
      ylabel style={above},
      colormap={blackwhite}{gray(0cm)=(0.4); gray(1cm)=(0.9)},
      view={0}{90}
    ]
      \addplot3[
        contour lua={
          levels={[[Z_LEVELS]]},
          labels=false,
        },
        samples=100,
        domain=[[X_MIN]]:[[X_MAX]],
        domain y=[[Y_MIN]]:[[Y_MAX]],
      ] {[[A]]*x^2 + [[B]]*y^2 + [[C]]*x*y + [[D]]*x + [[E]]*y + [[F]]};

      [[CYCLIC_PATH_ITER]]
        
    \end{axis}
  \end{tikzpicture}
  \caption{Итерация №[[K]] метода циклического покоординатного спуска.}
  \end{figure}

  \item Проверка критерия останова:
  \[
    [[BREAKPOINT_CHECK]]
  \]
\end{enumerate}